\maketitle

\begin{abstract}
\todo[inline]{Abstract --- written last, once the application numbers
are verified. Will state: (a) a delay-embedding local-Gaussian
conditional-kernel predictor instantiated on Ontario IESO demand;
(b) the out-of-sample backtest; (c) the directly-verified
infeasibility of a retrospective IESO head-to-head and the registered
forward experiment built in response; (d) the structural caveat that
\(\Pij\) departs from \(\kQ\) in both location and shape, with the
recoverable loss of each quantified and found small on the full
delivery-day population.}
\end{abstract}

\section{Introduction}
\label{sec:intro}

This paper frames short-horizon Ontario electricity-demand forecasting
as a question about the \emph{conditional law of a reconstructed
dynamical state}: given a delay embedding of recent demand, what is
the distribution of the next state? A local-Gaussian
conditional-kernel estimator of that law is instantiated on IESO data
and evaluated as a \emph{prediction experiment} --- with an
out-of-sample backtest, a registered forward experiment, and an
accounting of comparability against IESO's own published forecasts.

The estimator targets the conditional-regularity kernel \(\kQ\)
determined by the data measure via Rokhlin disintegration
(Section~\ref{sec:method}).

\todo[inline]{Prior-work context, for the author group to supply; this
draft makes no comparative-novelty claims.}

\section{Theory and method}
\label{sec:method}

\subsection{Construction}
\label{sec:construction}

\paragraph{Embedding and library.}
Let \(s_t\) be a scalar observable (here, de-seasonalised hourly
demand; Section~\ref{sec:pipeline}) sampled on a regular grid
\(t \in \mathbb{Z}\). The delay embedding into \(\Rb^{d}\) is
\begin{equation}
\label{eq:embed}
x_t \;=\; (s_t,\; s_{t-1},\; \ldots,\; s_{t-(d-1)}) \,.
\end{equation}
Per day-type \(\mathrm{dt} \in \{\mathrm{weekday},
\mathrm{saturday}, \mathrm{sunday}\}\) the embedding dimension
\(d = d(\mathrm{dt})\) is selected on pre-cutoff library data by the
elbow rule on the one-step in-sample prediction-skill curve
\(\rho \mapsto d\), with the \emph{day-type seam} excluded --- here
defined as the 07:00 hour at which the active day-type changes; a
one-step pair \((x_t, x_{t+1})\) whose target time \(t+1\) falls on
that hour crosses to a different day-type's library and is dropped
from the fit (see Section~\ref{sec:pipeline}). The pre-cutoff library
is
\begin{equation}
\label{eq:library}
\mathcal{L} \;=\; \{\,(x_t, x_{t+1}) : t \le t_{\mathrm{cut}},\;
t \text{ on the day-type } \mathrm{dt},\;
t+1 \notin \text{day-type seam}\,\} \,,
\end{equation}
with \(t_{\mathrm{cut}}\) the model-specification cutoff.

\paragraph{Local linear fit.}
For a query state \(x \in \Rb^{d}\), form the kernel weights
\begin{equation}
\label{eq:weights}
w_i(x;\theta) \;=\; (2\pi)^{-d/2}\,\theta^{-d}\,
\exp\!\left(-\frac{\lVert x_i - x \rVert^{2}}{2\theta^{2}}\right),
\qquad (x_i, y_i) \in \mathcal{L} \,,
\end{equation}
with bandwidth \(\theta > 0\) and \(y_i \equiv x_{t_i + 1}\). The
weighted least-squares local linear one-step map
\begin{equation}
\label{eq:drift}
x_{\text{next}} \;\approx\; x \cdot \Cj \qquad
\text{(row convention; \(j\) indexes the query state)}
\end{equation}
solves
\begin{equation}
\label{eq:wls}
\Cj \;=\; \arg\min_{C \in \Rb^{d \times d}}
\sum_{i \in \mathcal{L}} w_i(x;\theta) \, \lVert y_i - x_i C \rVert^{2}
\;=\; \bigl(X^{\top} W X\bigr)^{-1} X^{\top} W Y \,,
\end{equation}
where \(X, Y\) stack the library pairs and
\(W = \operatorname{diag}(w_i)\). The plain (unweighted)
residual covariance defines the diffusion,
\begin{equation}
\label{eq:diff}
\Sj \;=\; \frac{1}{|\mathcal{L}|}
\sum_{i \in \mathcal{L}} \bigl(r_i - \bar r\bigr)
\bigl(r_i - \bar r\bigr)^{\!\top},
\qquad r_i := y_i - x_i \Cj,
\quad \bar r := \tfrac{1}{|\mathcal{L}|}\sum_i r_i \,.
\end{equation}
The pair defines the state-dependent Gaussian Markov kernel
\begin{equation}
\label{eq:kernel}
\Pij(x, \cdot) \;=\; \Nrm\!\bigl(x \Cj,\; \Sj\bigr) \,.
\end{equation}

\paragraph{Why the plain (unweighted) covariance.}
The kernel-weighted residual covariance
\(\Sj^{w} = \sum_i w_i\,(r_i - \bar r)(r_i - \bar r)^{\!\top}
/ \sum_i w_i\) collapses toward \(\approx w^{2}\!\cdot Q\) for the
normalised Gaussian kernel of Eq.~\ref{eq:weights}; the plain
covariance is the unbiased local second moment, free of that
collapse.

\paragraph{Bandwidth: locality of \(\Cj\).}
Can the bandwidth \(\theta\) be chosen by data so that \(\Cj\) is
genuinely local? Two rules were examined. The Goldenshluger--Lepski
criterion in the supporting library predicts a bias--variance
interior optimum; for the normalised Gaussian kernel of
Eq.~\ref{eq:weights} its score is monotone in \(\theta\) with no
interior optimum, so it returns only the grid ceiling. True
leave-one-out cross-validation on one-step squared error,
\begin{equation}
\label{eq:loocv}
\mathcal{R}_{\mathrm{LOO}}(\theta)
\;=\; \frac{1}{|\mathcal{L}|}
\sum_{i \in \mathcal{L}}
\bigl\lVert y_i - x_i \Cj^{(-i)}(\theta) \bigr\rVert^{2},
\end{equation}
with \(\Cj^{(-i)}(\theta)\) the local fit at \(x_i\) excluding
\((x_i, y_i)\), does have an interior argmin, but under a
bootstrap-SE procedure (\(B = 10^{3}\) library resamples) it is not
statistically separated from any other grid point. The same
non-discrimination reproduces on three controlled systems (a linear
VAR(1), the logistic map, a regime-switching AR) across horizons
\(h \in \{1, 2, 4, 8, 16, 24\}\): it is a property of the criterion,
not of Ontario data. At the \(d = 2\) z-score-lag embedding, \(\Cj\)
is therefore effectively a single global linear operator.

With \(\Cj\) effectively global and linear, the non-Dirac content of
the predictive law is carried by \(\Sj\) and its departure from
Gaussianity (Section~\ref{sec:caveat}). The available lever is the
\emph{state representation} (richer embedding, explicit
phase/derivative coordinate), not the bandwidth rule.

\paragraph{The conditional law \(\kQ\) the estimator targets.}
On a probability space \((\Omega, \mathcal{F}, P)\) with two
observations \(Q, F\), the Rokhlin disintegration
theorem~\cite{Rokhlin1952,Bogachev2007} produces a Markov kernel
\begin{equation}
\label{eq:kappaQ}
\kQ(q, \cdot) \;=\; P(F \in \cdot \mid Q = q),
\end{equation}
determined by the measure --- not posited as an additional dynamical
law. Under temporal indexing of the observations, finite-dimensional
distributions are coherent, so the kernels compose: the
Chapman--Kolmogorov equation is a consequence, not a hypothesis, and
yields a Markov semigroup \(\{\Pi_t\}\)~\cite{Eisner2015}. The
companion Koopman--Perron duality
\begin{equation}
\label{eq:KP}
\int K_t\, g\, \dd\mu \;=\; \int g\, \dd(P_t^{*}\mu)
\end{equation}
relates the Koopman operator \(K_t g = g \circ T_t\) (when the
dynamics is deterministic, recovering the classical Koopman operator
when the kernels are Dirac~\cite{LasotaMackey1994}) to its
Perron--Frobenius adjoint \(P_t^{*}\) on densities. Three layers are
useful to keep notationally distinct: \(\kQ\) (the disintegration
kernel; a conditional law), \(\{\Pi_t\}\) (the time-indexed Markov
semigroup), and \(K_t\) (the operator-layer object).

\paragraph{The estimator's relationship to \(\kQ\).}
The object Section~\ref{sec:construction} constructs is a
finite-sample, parametric, single-step, Gaussian-form estimator of
\(\kQ\) at the delay-embedding layer, from which the operator-layer
objects \(K_{\Delta} / P_{\Delta}^{*}\) are recoverable in the
non-Dirac regime (\(\Sj \neq 0\)).
The \(\kQ\) correspondence holds at the embedding layer; \emph{locality}
of that estimator is not empirically operative on Ontario
(Section~\ref{sec:construction}).
\(\Cj\) is a linearisation of the conditional-mean factor (intended
local, effectively global here); \(\Sj\) is the second moment of
\(\kQ\)'s Gaussian projection; \(\Sj \to 0\) reproduces the
classical-Koopman (deterministic) collapse. Three qualifiers bound
this correspondence:

\begin{enumerate}
  \item \textbf{Register.} \(\kQ\) is determined by the measure; this
  work \emph{constructs} a Gaussian proxy and fits it from finite
  data --- ``estimator of,'' not ``instance of.''

  \item \textbf{Imposed Gaussianity.} \(\Pij = \Nrm(x \Cj, \Sj)\)
  coincides with \(\kQ\) only where the local conditional law is
  itself Gaussian. Section~\ref{sec:caveat} quantifies the gap as
  real but small in recoverable forecast loss.

  \item \textbf{Single step.} \(\Pij\) is the single-step slice
  \(\Pidelta\); the semigroup \(\{\Pi_t\}\) is neither composed nor
  verified here. (A multi-step predictive-variance propagation via a
  shift-aware state Jacobian is built and gate-validated as the
  instrument for Section~\ref{sec:caveat}'s decomposition; it is not
  the semigroup.)
\end{enumerate}

\section{Structural caveat: a mis-centred conditional mean and a
Gaussian-proxy diffusion}
\label{sec:caveat}

The two estimator factors \(\Cj\) and \(\Sj\) each carry a quantified
limitation. A decomposition on the 2021--22 development-set window,
with predictive variance propagated over the full delivery-day
population, finds \textbf{neither effect large enough to act on}.

\paragraph{\(\Sj\) is rank-1 by construction.} For a single-variable
delay embedding, coordinates \(2, \ldots, d\) of the one-step image
are deterministic shifts of the input coordinates
(\(Y_{[:,\,1:]} = X_{[:,\,:-1]}\) exactly). The only stochastic
content of the one-step map is the scalar coordinate-0 innovation;
therefore \(\Sj\) is rank-1 by construction, independent of data ---
the multivariate diffusion is structurally a single scalar innovation
variance embedded in \(\Rb^{d \times d}\).

\paragraph{That scalar innovation is non-Gaussian.} The
non-Gaussianity question is therefore \emph{univariate}. The
estimator, validated against synthetic ground truth
(Section~\ref{sec:synth}), gives a heavy-tailed one-step innovation.

\paragraph{The conditional mean is phase-mis-centred.}
The conditional-mean factor is not faithful either. The embedding
gives the local linear map no phase information --- it cannot tell a
state rising through a demand level from one falling through it ---
so it centres between rising and falling regimes and carries a
structured residual bias. The signature is a correlation between the
signed \(z\)-space error and the demand derivative \(\dd z/\dd t\):
measured, it is \(\mathrm{corr} = -0.246\) (bootstrap 95\% CI
\([-0.283,\, -0.207]\), \(n = 3167\)), monotone across
rising/falling phase. \(\Pij\)'s location, not only its shape,
departs from \(\kQ\).

\paragraph{Decomposition.}
\label{sec:caveat-decomposition}
A diagnostic on the 2021--22 development-set window bounds the
recoverable loss of each factor by oracle counterfactuals on the
multi-step forecast log-loss \(\Lbar\), over all delivery days, with
predictive variance propagated across the embedding-dimension change.
The recoverable-loss factors are
\begin{equation}
\label{eq:DA-DB}
\DA \;:=\; \frac{\Lbar - \Lbar^{(A)}}{\Lbar},
\qquad
\DB \;:=\; \frac{\Lbar - \Lbar^{(B)}}{\Lbar},
\end{equation}
where \(\Lbar^{(A)}\) is the log-loss under the oracle that removes
the phase-conditional mean bias, and \(\Lbar^{(B)}\) is the loss
under the oracle that replaces the Gaussian innovation law with a
Student-\(t\) fitted on the residuals (shape \(\nuhat\)).
On the full delivery-day population the recoverable loss is small in
both directions: \(\DA = 0.048\) (95\% day-bootstrap CI
\([0.007,\,0.086]\)) and \(\DB = 0.023\) (CI
\([-0.000,\,0.048]\), including zero), at \(\nuhat = 6.5\). Neither
effect is large or cleanly established enough to act on; restricting
to the single-dimension subpopulation (\(\approx 58\%\) of days)
gives larger point estimates (\(\DA \approx 0.061\),
\(\DB \approx 0.091\)) but overlapping intervals, the same
qualitative reading. The diffusion-shape diagnostic is run on
standardised residuals of the multi-step forecast under
propagation-aware variance; on the same data the propagation softens
the apparent kurtosis roughly four-fold
(\(\mathrm{kurt}_{\mathrm{exc}}: 6.1 \to 2.4\)).

\paragraph{The one-step conditional law remains heavy-tailed.} The
propagation softening is a property of the multi-step forecast
object, not of the underlying conditional law: the per-anchor
library residual has \(\mathrm{kurt}_{\mathrm{exc}} \approx
24\)--\(33\) (one-step, local, no propagation), and that figure is
preprocessing-independent. Section~\ref{sec:qualifier-verdict}
quantifies this one-step departure directly and examines what can be
done about it.

\section{Application and results}
\label{sec:application}

\subsection{Ontario pipeline}
\label{sec:pipeline}

The instantiation is a configuration-driven pipeline; all paths and
parameters resolve from a single \texttt{config/pipeline.yaml}
relative to the project root (no working-directory dependence).
Stages: \textbf{acquire} IESO public \texttt{PUB\_Demand};
\textbf{de-seasonalise} (Section~\ref{sec:zscore}) hourly demand to a
stationary \(z\)-score; \textbf{day-type tag}
\(\{\mathrm{weekday}, \mathrm{saturday}, \mathrm{sunday}\}\) using a
configurable 07:00 day-anchor offset (transitions across the day-type
seam are excluded as in Section~\ref{sec:construction} --- without
this \(\Sj\) is catastrophically ill-conditioned, so the estimand is
effectively forced to be \emph{intra-day}); \textbf{embedding
dimension} selected per day-type by an elbow rule on the
prediction-skill-vs-dimension curve; \textbf{per-anchor fit} of the
Section~\ref{sec:method} estimator.

\subsection{The \(z\)-score transform is a modelling choice, not a
preprocessing detail}
\label{sec:zscore}

De-seasonalisation is
\begin{equation}
\label{eq:zscore}
z \;=\; \frac{D - \mu_{m,h}}{\sigma_{m,h}},
\end{equation}
with \(\mu, \sigma\) a month \(\times\) hour-of-day climatology
estimated strictly from pre-cutoff data. This single climatology
object enters \emph{all four} estimator layers: the embedded state,
the drift fit, the diffusion (\(\Sj\) is the covariance of
\(z\)-residuals), and the demand-space round-trip. It is therefore an
implicit seasonal model coupled into \(\kQ\), not a benign
preprocessing step, and its specification is a modelling choice worth
examining directly.

Two questions were investigated on the development-set window. The
first: is the phase-asymmetry structure of Section~\ref{sec:caveat}
an artefact of de-seasonalisation? Measured in raw-demand units so it
is comparable across specifications, the phase-asymmetry correlation
is \(\mathrm{corr}(\dd D/\dd t,\, e) = -0.247,\, -0.250,\, -0.217\)
under the global, month\,\(\times\)\,hour, and no-transform
specifications respectively --- near-identical, and present already
in raw demand. The phase-mis-centring is a property of the demand
dynamics, not a transform artefact.

The second: does de-seasonalisation earn its place on forecast
accuracy? Re-fitting the pipeline end-to-end on raw demand and on a
scale-only variant (\(z = D/\sigma_{m,h}\), no level removal), the
full month\,\(\times\)\,hour transform is decisively better ---
\(\MAPE\) on the development window degrades from \(3.86\%\) (full
transform) to \(11.98\%\) (scale-only) to \(13.57\%\) (raw demand).
The interval-coverage breakdown locates the mechanism: the raw-demand
pipeline delivers uniform-width intervals that over-cover overnight
(demand small and stable) and under-cover at the midday peak ---
the signature of a missing demand-space scale step that the
\(\sigma_{m,h}\) factor supplies. The de-seasonalisation transform
is load-bearing for forecast accuracy; both its level and scale
sub-steps contribute, and neither is removable without a substantial
accuracy cost. The month\,\(\times\)\,hour climatology is therefore
retained as the production choice; Section~\ref{sec:backtest}
compares it at production scale against a continuous-basis
alternative.

\subsection{Levers against the heavy-tailed conditional law}
\label{sec:qualifier-verdict}

Given the heavy-tailed one-step innovation of
Section~\ref{sec:caveat}, what can be done within the past-demand
information set? The \(\kQ\) construction admits four candidate
design changes that stay within the lagged-\(z\) embedding: a
Student-\(t\) local kernel in place of the Gaussian; state
augmentation within the embedding; per-horizon direct fitting; and
the combined state-augmentation-plus-kernel variant. Each was tested
on the development set, and the result is read against an independent
measurement of how much non-Gaussian mass is actually present.

\emph{Contamination magnitude.} A Patra--Sen
consistent two-component-mixture estimator~\cite{PatraSen2016},
with the predicted-Gaussian as the known component, applied to
pooled standardised per-anchor library residuals from the
production estimation path, places at least \(49\)--\(74\%\) of
the per-anchor library mass in the non-Gaussian mixture component
across day-types and preprocessings (\(\hat\alpha_L^{95\%}\), a
conservative lower bound by their Lemma~4). A pipeline-aware
Gaussian-innovation null through the identical protocol gives
\(\hat\alpha_L^{95\%} = 0.002\); the maximum value any single
heavy-tailed AR(1) reproduces, across both population-kurt-matched
and sample-kurt-matched checks, is \(0.26\). The contamination is
mixture-structured, not a small heavy-tail perturbation, and is
preprocessing-independent.

\emph{Empirical result.} On the development set (\(n = 132\) days),
only the Student-\(t\) kernel delivers a clean log-loss improvement
on routine-day data (\(\approx 1.8\%\) at mid horizons, CI excludes
zero); state augmentation, per-horizon direct fitting, and the
combined variant are null. The improvement is modest --- far below
the contamination mass just measured. The consistent reading of
large contamination alongside modest-to-null levers is that the
contamination is structurally orthogonal to the lagged-\(z\) state on
which all four levers operate. Reducing it requires exogenous
information --- temperature, weather regime, calendar event flags ---
not further variation within the past-demand information set.

\subsection{Out-of-sample backtest}
\label{sec:backtest}

The iterated one-step predictor is backtested on Ontario ground truth
for complete delivery days strictly after the model-specification
cutoff. Leakage guard: fits and \(z\)-score climatology use only
\(\le\)-cutoff data; post-cutoff actuals enter only as query history
and scoring truth. With forecast \(\hat{y}_{d,h}\) and actual
\(y_{d,h}\) over delivery days \(d \in \mathcal{D}\) and horizons
\(h \in \{1, \ldots, 24\}\),
\begin{equation}
\label{eq:mae-mape}
\MAE \;=\; \frac{1}{|\mathcal{D}|\cdot 24}
\sum_{d \in \mathcal{D}} \sum_{h=1}^{24}
\bigl| \hat{y}_{d,h} - y_{d,h} \bigr|,
\qquad
\MAPE \;=\; \frac{1}{|\mathcal{D}|\cdot 24}
\sum_{d \in \mathcal{D}} \sum_{h=1}^{24}
\frac{\bigl| \hat{y}_{d,h} - y_{d,h} \bigr|}{y_{d,h}} \,.
\end{equation}
In a day-ahead forecast, horizon \(h\) is collinear with hour-of-day,
so the error structure is \emph{diurnal} (deep night easiest;
dawn/dusk ramps hardest), not pure horizon compounding. The recorded
backtest gives \(\MAE = 786~\mathrm{MW}\),
\(\MAPE = 4.78\%\) over \(|\mathcal{D}| = 500\) out-of-sample
days, \(\approx 1.5\times\) better than seasonal persistence at
every hour. The number is bound to the recorded artifact by its
content hash;\footnote{Body sha256 \texttt{e405a121\(\ldots\)};
recorded under frozen-spec hash \texttt{f0dd9c87\(\ldots\)}
(\(\mu_{m,h}\) climatology). The per-horizon table appears in
Appendix~\ref{app:backtest}.}
interval coverage carries the Section~\ref{sec:caveat}
Gaussian-proxy caveat: a Gaussian interval around a heavy-tailed
innovation under-covers.

\paragraph{Climatology comparison.} The \((\mu_{m,h}, \sigma_{m,h})\)
lookup is piecewise-constant and jumps discretely at month boundaries
(twelve delivery days per year). A continuous Fourier climatology ---
a basis in (day-of-year, hour-of-day) with
\(K_{\mathrm{year}} = K_{\mathrm{day}} = 8\), no interaction terms ---
was re-fit and run on the same \(500\)-day post-cutoff span as an
alternative. The two backtests compare as follows:
\begin{center}\small
\begin{tabular}{l r r r}
\toprule
metric & \(\mu_{m,h}\) lookup & Fourier (\(K_y=K_d=8\)) & \(\Delta\) \\
\midrule
overall \(\MAE\)            & \(786\)~MW   & \(904\)~MW   & \(+118\)~MW \\
overall \(\MAPE\)           & \(4.78\%\)   & \(5.42\%\)   & \(+0.64\)~pp \\
persistence-168h \(\MAE\)   & \(1138\)~MW  & \(1138\)~MW  & --- \\
interval coverage           & \(20.8\%\)   & \(21.8\%\)   & \(+1.0\)~pp \\
\bottomrule
\end{tabular}
\end{center}
The \(\mu_{m,h}\) lookup wins overall by \(\sim 15\%\) in \(\MAE\).
A diurnal breakdown locates the difference: Fourier improves slightly
on the dawn ramp (\(h = 7\)--\(9\)) but loses substantially across
the daytime peak (\(h = 11\)--\(16\), \(+161\) to \(+361\) MW), where
demand --- and thus absolute error --- is highest. The mechanism is
that the lookup's \(288\) independent cells carry the day-of-year
\(\times\) hour-of-day interaction explicitly (the seasonal peak-load
shape varies by month), which a Fourier basis without cross-terms
cannot represent. The month-boundary jumps are real but of small
operational consequence: the \(\sim 12\) days per year on which
Fourier improves do not outweigh the \(\sim 365\) on which it loses
seasonal-by-hour resolution.\footnote{A Fourier basis \emph{with}
cross-interaction terms (\(\sim 289\) parameters, matching the
lookup's cell count) is a natural variant, out of scope here.}

\subsection{Registered forward prediction experiment}
\label{sec:registered}

A fair retrospective IESO comparison is not recoverable
(Section~\ref{sec:ieso}), so the comparison is run forward. The model
specification (estimator identity, resolved hyperparameters, data
cutoff \(t_{\mathrm{cut}}\), code commit, and a content hash of the
\(\le t_{\mathrm{cut}}\) actuals) is committed before any forecast
is issued, and the same content hash is re-checked at every
prediction-time load so the recorded specification cannot be silently
edited. The predictor is local/lazy (per-anchor fit at prediction
time from \(\le t_{\mathrm{cut}}\) library data), so the committed
specification pins the function and its inputs; a cutoff guard blocks
any post-\(t_{\mathrm{cut}}\) data from entering a fit.
Forecasts issue forward; scraped actuals settle into an append-only
ledger (IESO restates recent demand --- revisions appended, never
overwritten); scoring runs against persistence baselines with IESO
published values recorded as differently-horizoned context.

\todo[inline]{Forward-experiment results accrue with calendar time;
report once a sufficient settled window exists.}

\subsection{The retrospective IESO head-to-head is infeasible from
public archives}
\label{sec:ieso}

Can the forecast be compared retrospectively against IESO's own
published forecast on historical settled days? A fair comparison
requires two properties of an IESO archive: (i) the values reported
must be on the Ontario-demand basis (matching the forecast target
of this paper), and (ii) the values reported must be the day-ahead
issue (matching its horizon). The simpler hypothesis is that some
public IESO archive supplies both. To test it, all 133 public IESO
report directories were enumerated and each checked against the two
properties.

The simpler hypothesis fails, via two independent blockers. The
Ontario-demand-basis archive (Adequacy3
\texttt{ForecastOntDemand}) retains, for days old enough to have
settled actuals, only the end-of-delivery-day version --- so its
horizon is wrong (property (ii) violated). The day-ahead-issued
archives (DATotals, PredispTotals) publish only market totals --- so
their basis is wrong (property (i) violated): ``Total Load'' carries
a definitional \(\approx +2485~\mathrm{MW}\) offset over Ontario
Demand (MarketQuantity enumeration: Total
Energy/Loss/Load/Dispatchable/10S/10N/30R; none is Ontario demand).
No archive satisfies both. The retrospective comparison is
infeasible from public data; hence the registered forward design
of Section~\ref{sec:registered}. (Stated for public archives as of
the recorded access date.)

\subsection{Validation against synthetic ground truth}
\label{sec:synth}

Three gates, summarised in Table~\ref{tab:gates}, all pass on
known-ground-truth synthetic systems.

\begin{table}[h]
\centering
\small
\caption{Synthetic-ground-truth validation gates. Recorded result
artifact body \(\text{sha256} = \text{d843}{\ldots}\).}
\label{tab:gates}
\begin{tabular}{l l p{0.42\linewidth}}
\toprule
Gate & System & Result \\
\midrule
Recovery
  & Known-\((A,Q)\) VAR(1)
  & drift rel.\ err.\ \(0.048\) (tol.\ \(0.10\)), diffusion
    \(0.026\) (tol.\ \(0.25\)) \\
Non-Gaussianity
  & Gaussian / \(t_5\) VAR(1)
  & excess kurt.\ \(-0.05\) / \(+4.72\); heavy tail flagged \\
Multi-step
  & \(\Pidelta^{k}\) vs.\ closed form
  & drift rel.\ err.\ \(0.05 \to 0.44\) over \(h = 1 \to 24\),
    within tolerance \\
\bottomrule
\end{tabular}
\end{table}

\noindent
The recovery gate confirms \(\Cj\) and \(\Sj\) are estimated within
tolerance on a system whose \((A, Q)\) are known; the
non-Gaussianity gate confirms the diagnostic reads a Gaussian
innovation as \(\approx\)Gaussian and a \(t_5\) innovation as
heavy-tailed; the multi-step gate shows the iterated map's drift
error growing with horizon but staying within the
horizon-dependent tolerance band through \(h = 24\). All three are
recorded as a single CLAIM-grade artifact (body sha256
\texttt{d8435415\(\ldots\)}).

\section{Limitations and reproducibility}
\label{sec:limits}

\begin{itemize}
  \item \textbf{Estimator limitations (Section~\ref{sec:caveat}).}
  \(\Pij\) departs from \(\kQ\) in location (\(\Cj\)
  phase-mis-centred) and shape (\(\Sj\) Gaussian proxy of a
  heavy-tailed innovation). The decomposition bounds the recoverable
  loss at \(\approx 4.8\%\) and \(\approx 2.3\%\) of development-set
  log-loss --- small in both directions, with the heavy-tail interval
  including zero. Multi-step propagation softens the apparent
  kurtosis \(\approx 4\times\); the one-step conditional law remains
  heavy-tailed at the library-residual level.

  \item \textbf{Within-\(z\) levers (Section~\ref{sec:qualifier-verdict}).}
  Of the four candidate design changes within the lagged-\(z\)
  embedding --- a Student-\(t\) kernel, state augmentation,
  per-horizon direct fitting, and the combined variant --- only the
  Student-\(t\) kernel delivers a clean log-loss improvement, and
  the effect is modest (\(\approx 1.8\%\) at mid horizons). State
  augmentation and the combined variant are null.

  \item \textbf{Single-step only.} The semigroup \(\{\Pi_t\}\) is not
  constructed; multi-step is iteration of \(\Pidelta\), validated for
  error growth. A semigroup treatment is open.

  \item \textbf{Exogenous information.} No weather/exogenous
  covariates. The Patra--Sen quantification places at least
  \(49\)--\(74\%\) of per-anchor library mass in a non-Gaussian
  mixture component that the within-\(z\) levers do not access;
  reducing the heavy-tailed departure requires exogenous information
  rather than further variation within the \(z\) information set.
\end{itemize}

\paragraph{Reproducibility.} Recorded results carry two content
hashes: an inputs hash (git commit, model-specification hash, library
versions, RNG seeds, declared inputs) and a body hash recomputed at
read time. The recording refuses to write from a dirty source tree.
Exploratory development-set material is marked as such and is not
cited as a result elsewhere in the paper.

\section{Conclusion}
\label{sec:conclusion}

The two ways \(\Pij\) departs from \(\kQ\) --- a phase-mis-centred
conditional mean and a Gaussian-proxy diffusion --- were each
quantified by an oracle decomposition, and the recoverable loss is
small in both directions on the full delivery-day population. Of the
four design changes available within the lagged-\(z\) embedding, only
the Student-\(t\) kernel delivers a clean log-loss improvement, and
it is modest (\(\approx 1.8\%\) at mid horizons). A Patra--Sen
contamination estimate places at least \(49\)--\(74\%\) of the
per-anchor library mass in a non-Gaussian mixture component ---
substantial mass the within-\(z\) levers do not reach. The consistent
reading is that the contamination is structurally orthogonal to the
lagged-\(z\) state: reducing the heavy-tailed departure of the
conditional law requires \emph{exogenous information} (temperature,
weather regime, calendar event flags), not further variation within
the past-demand information set. The estimator's evaluation is run
forward as a registered prediction experiment, the retrospective
IESO head-to-head being infeasible from public archives.

\appendix

\section{Estimator detail}
\label{app:estimator}

\todo[inline]{WLS weighting, the LOO-CV \(\theta\) rule, the
GL-degeneracy note in full.}

\section{Per-horizon backtest table}
\label{app:backtest}

Table~\ref{tab:backtest-horizons} reports per-horizon backtest error
on the out-of-sample delivery days
(\(|\mathcal{D}| = 500\) days, horizons \(h = 1, \ldots, 24\)).
\(\MAE\) and \(\MAPE\) are defined in Eq.~\ref{eq:mae-mape}; the
persistence baseline is the seasonal-persistence value at the same
hour-of-day from one week prior.

\begin{table}[h]
\centering
\caption{Per-horizon backtest error and ratio to the
seasonal-persistence baseline.}
\label{tab:backtest-horizons}
\begin{tabular}{c c c c}
\toprule
Horizon \(h\) & \(\MAE_h\) (MW) & \(\MAPE_h\) (\%) & \(\MAE_h /\MAE_h^{\mathrm{persist}}\) \\
\midrule
\(1\)   & \(161\)  & \(1.08\) & \(0.17\) \\
\(2\)   & \(215\)  & \(1.48\) & \(0.23\) \\
\(3\)   & \(241\)  & \(1.67\) & \(0.26\) \\
\(4\)   & \(288\)  & \(2.00\) & \(0.32\) \\
\(5\)   & \(370\)  & \(2.54\) & \(0.40\) \\
\(6\)   & \(563\)  & \(3.80\) & \(0.59\) \\
\(7\)   & \(839\)  & \(5.43\) & \(0.83\) \\
\(8\)   & \(1011\) & \(6.25\) & \(0.96\) \\
\(9\)   & \(1056\) & \(6.43\) & \(0.96\) \\
\(10\)  & \(1130\) & \(6.92\) & \(0.94\) \\
\(11\)  & \(1190\) & \(7.27\) & \(0.92\) \\
\(12\)  & \(1182\) & \(7.23\) & \(0.87\) \\
\(13\)  & \(1154\) & \(7.10\) & \(0.82\) \\
\(14\)  & \(1105\) & \(6.82\) & \(0.78\) \\
\(15\)  & \(1037\) & \(6.36\) & \(0.74\) \\
\(16\)  & \(959\)  & \(5.72\) & \(0.71\) \\
\(17\)  & \(891\)  & \(5.08\) & \(0.71\) \\
\(18\)  & \(886\)  & \(4.89\) & \(0.74\) \\
\(19\)  & \(908\)  & \(4.95\) & \(0.78\) \\
\(20\)  & \(893\)  & \(4.93\) & \(0.78\) \\
\(21\)  & \(857\)  & \(4.86\) & \(0.75\) \\
\(22\)  & \(745\)  & \(4.40\) & \(0.67\) \\
\(23\)  & \(612\)  & \(3.75\) & \(0.58\) \\
\(24\)  & \(579\)  & \(3.67\) & \(0.57\) \\
\midrule
overall & \(786\) & \(4.78\) & \(0.69\) \\
\bottomrule
\end{tabular}
\end{table}

The error is diurnal, not horizon-compounding: \(\MAE_h\) is lowest
overnight (\(h \le 5\)) and peaks across the midday demand plateau
(\(h = 11\)--\(13\)). The ratio to seasonal persistence is best
overnight (\(0.17\) at \(h = 1\)) and weakest across the dawn ramp
(\(h = 8\)--\(9\), ratio \(0.96\)) --- the dawn ramp is the hardest
part of the day for both the predictor and the persistence
baseline.

% ---------- References ----------
\ifstandalone
\bibliographystyle{plain}
\bibliography{references}
\fi
